\documentclass[a4paper,12pt]{article}
\usepackage[utf8]{inputenc}
\usepackage{enumitem}
\usepackage{geometry}
\geometry{margin=2.5cm}
\usepackage{titlesec}
\titleformat{\section}{\large\bfseries}{\thesection}{1em}{}

\title{Plan de Ejecución Procedural en Unity (MVP)}
\author{Donovan}
\date{\today}

\begin{document}

\maketitle

\section{Objetivo}
Diseñar un sistema de generación procedural en Unity 2D que:
\begin{itemize}
    \item Genere mapas únicos mediante una estrategia de \textbf{grid}.
    \item Instancie bloques de escenario a partir de un \textbf{banco de prefabs}.
    \item Use una \textbf{semilla} para reproducibilidad.
    \item Controle colisiones entre objetos instanciados.
    \item Limite intentos de colocación mediante un \textbf{máximo de tries}.
\end{itemize}

\section{Estructura de Objetos}
\begin{itemize}
    \item \textbf{BlockBankSO (ScriptableObject)}  
    Contiene la lista de prefabs disponibles por nivel.  
    \begin{itemize}
        \item \texttt{List<GameObject> blockPrefabs}
        \item \texttt{GetRandomBlock(seed)}: devuelve un prefab aleatorio según semilla.
    \end{itemize}

    \item \textbf{MapGenerator (MonoBehaviour)}  
    Responsable de recorrer el grid y generar el mapa.  
    \begin{itemize}
        \item Parámetros: \texttt{width}, \texttt{height}, \texttt{cellSize}, \texttt{seed}.
        \item Lógica: recorre el grid, pide bloques al \texttt{BlockBankSO}, instancia en posiciones discretas.
        \item Aplica rotación aleatoria.
    \end{itemize}

    \item \textbf{BlockContent (MonoBehaviour)}  
    Script en cada prefab que genera entidades internas.  
    \begin{itemize}
        \item \texttt{GameObject[] entities}
        \item \texttt{spawnCount}
        \item Método: spawnea entidades en posiciones aleatorias dentro del bloque.
    \end{itemize}

    \item \textbf{Spawner (MonoBehaviour)}  
    Controla la instanciación de entidades con chequeo de colisiones.  
    \begin{itemize}
        \item \texttt{List<GameObject> spawnedObjects}
        \item \texttt{maxTries}
        \item Método: valida que el punto candidato no caiga dentro del collider de otro objeto.
    \end{itemize}
\end{itemize}

\section{Flujo Ideal}
\begin{enumerate}
    \item \textbf{Inicialización:}
    \begin{itemize}
        \item Se carga el \texttt{BlockBankSO} correspondiente al nivel.
        \item Se define la semilla para el generador aleatorio.
    \end{itemize}

    \item \textbf{Generación del Grid (MapGenerator):}
    \begin{itemize}
        \item Recorre cada celda del grid.
        \item Solicita un prefab aleatorio al \texttt{BlockBankSO}.
        \item Instancia el prefab en la posición correspondiente con rotación aleatoria.
    \end{itemize}

    \item \textbf{Generación Interna del Bloque (BlockContent):}
    \begin{itemize}
        \item Cada bloque instanciado ejecuta su script.
        \item Se generan entidades internas (árboles, rocas, enemigos) en posiciones aleatorias.
    \end{itemize}

    \item \textbf{Control de Colisiones (Spawner):}
    \begin{itemize}
        \item Antes de instanciar una entidad, se valida contra la lista de objetos ya instanciados.
        \item Se comprueba que el punto candidato no esté dentro del collider de otro objeto.
        \item Si hay colisión, se genera un nuevo punto y se repite hasta \texttt{maxTries}.
        \item Si no se encuentra espacio válido, se aborta la instanciación.
    \end{itemize}

    \item \textbf{Resultado:}
    \begin{itemize}
        \item Se obtiene un mapa único, reproducible por semilla.
        \item Cada bloque contiene entidades distribuidas sin colisiones.
    \end{itemize}
\end{enumerate}

\section{Consideraciones de Rendimiento}
\begin{itemize}
    \item Prefabs optimizan la instanciación y permiten reutilización.
    \item Uso de Tilemaps mejora eficiencia en escenarios 2D.
    \item El control de colisiones con \texttt{maxTries} evita bucles infinitos.
\end{itemize}

\section{Conclusión}
El sistema propuesto combina:
\begin{itemize}
    \item \textbf{Grid procedural} para estructura base.
    \item \textbf{Prefabs} para bloques y entidades.
    \item \textbf{ScriptableObject} para bancos de datos.
    \item \textbf{Colisión + máximo de tries} para control de instanciación.
\end{itemize}

\end{document}
